\subsection{Baseline details}
\label{app:baselines}

All baselines operate under the same budget as our method
(Def.~\ref{def:budget}): unrestricted source interaction, the same $20$ target
demonstrations, and at most $100$k target transitions. None of them sees target
reward shaping, the target physics parameters, or a target expert. Every arm is
evaluated in the same target environment with the same success predicate and the
same number of episodes.

Two budget conventions are worth stating, because they are easy to get wrong and
they change the numbers. First, the warm-up episodes that each online method
runs before its first gradient step \emph{are} counted against the $100$k; we
verified this by checking that the step counter includes them, rather than
assuming it. Second, the $20$ demonstrations are \emph{not} counted, for any
method, since every method gets the same ones. Warm-up rollouts are collected
with the arm's own frozen prior acting deterministically, so a method that warm
starts its critic does not thereby get better warm-up data.

\paragraph{IBRL.}
Our model-free reference, and the strongest of the four.
Following \citet{hu2023ibrl}, a behaviour-cloning prior $\mu$ is obtained by
fine-tuning the source expert on the target demonstrations and then frozen. An
actor $\pi_\phi$ and a twin critic $Q_\psi$ are trained online with TD3-style
updates, and at every step the executed action is the critic's argmax over the
two candidates $\{\pi_\phi(s), \mu(s)\}$ (Eq.~\ref{eq:ibrl}). The same argmax is
used inside the Bellman backup, so the prior bootstraps the value function
rather than only regularising the actor. The demonstrations are kept in a frozen
half of the replay buffer, so every minibatch is half demonstration data and
half online data throughout training.

We report three variants, which differ in exactly one thing each:
\begin{itemize}
\item \textbf{IBRL} --- the published method: the critic is initialised from
      scratch.
\item \textbf{IBRL + source critic} --- the critic is initialised from the
      source expert's critic $Q_s$ instead. This is the only change; the prior,
      the data, and every hyperparameter are identical. We add this arm because
      the model-based methods already carry the source value (their value head
      is distilled from $Q_s$), so comparing them against a from-scratch critic
      credits the world model with something the model-free method was simply
      never given.
\item \textbf{IBRL + source critic + margin} --- additionally requires the
      learned actor to beat the prior by a margin $m$ before it is executed
      (Sec.~\ref{sec:two-conditions}). Included to show the margin is not
      specific to the model-based setting.
\end{itemize}

\paragraph{Dyna-IBRL.}
The model-based counterpart to IBRL, and the arm that isolates \emph{how} a
world model is used. It shares everything with our method's stages~1 and 2 ---
the same distilled world model, the same frozen prior from the same
tethered warm-up --- and differs only in what happens during stage~3.
Instead of planning with the model at decision time, it uses the model as a data
generator: each round it branches imagined rollouts from real latents, adds them
to a synthetic buffer, and trains an actor and a $Q(z,a)$ head on a mixture of
real and imagined transitions. At deployment the actor runs alone, with no
planner.

Three implementation points matter for a fair comparison. The $Q$ head does not
exist in the distilled model, which provides $\hat r$ and $\hat V$ but not
$Q(z,a)$; we therefore add one and initialise it by regression onto
$\hat r(z,a) + \gamma \hat V(f(z,a))$, so that at step zero the critic
\emph{is} the source value and every subsequent change is due to target data.
Imagination is carried out entirely in latent space rather than decoding to
states and re-encoding, because that round trip is not the identity and its
error is of the same order as the one-step latent motion. And the imagined
rollouts are collected with the same IBRL argmax that acts in the environment,
not with the actor alone, so the critic is trained on the distribution it will
actually be evaluated under.

We report Dyna-IBRL with and without the prior-preference margin, and with and
without the dynamics ensemble, since Sec.~\ref{sec:two-conditions} shows these
two are not independent.

\paragraph{RLPD.}
An off-policy method designed to learn efficiently from a mixture of offline and
online data \citep{ball2023rlpd}, included as a representative of the
``no prior policy, but plenty of offline data'' regime. Following the original
recipe we use symmetric sampling --- every minibatch is half offline and half
online --- together with layer normalisation in the critic, a critic ensemble
with random subset selection, and a high update-to-data ratio.

The offline buffer is the union of a large set of source rollouts and the $20$
target demonstrations, merged into a single prior dataset. This gives RLPD
strictly more offline data than any other arm receives. We do so deliberately:
RLPD has no mechanism for using a source \emph{policy}, so giving it the source
\emph{data} is the closest equivalent, and it prevents the comparison from
reducing to ``the method without a prior does worse''.

\paragraph{IQL.}
An offline-RL method \citep{kostrikov2022iql} used here as a warm start followed
by online fine-tuning, which is the setting in which it is usually deployed for
sim-to-real. We first run IQL offline on source rollouts plus the target
demonstrations: the value function is fit by expectile regression, which avoids
querying the critic on actions absent from the data, and the policy is extracted
by advantage-weighted regression. The resulting actor and critic then initialise
an online run with the same budget and environment as the other arms.

One deviation from the standard recipe is worth recording. On this data
advantage-weighted regression degenerates to plain behaviour cloning, because
the spread of advantages across a minibatch is far smaller than the spread of
$Q$ values, so the exponential weights are nearly uniform. We therefore use
reward-weighted regression for policy extraction instead. This is a change that
\emph{helps} the baseline; we report it because the failure mode is invisible in
the training curves and would otherwise look like a tuning problem.

\begin{table}[h]
\centering
\caption{Baseline settings. The upper block is shared by every online arm and is
listed once; the lower blocks give what is specific to each method. A dash means
the method has no such component.}
\label{tab:app-baselines}
\small
\begin{tabular}{ll}
\toprule
\multicolumn{2}{l}{\emph{Shared by all online arms}} \\
\midrule
target budget                     & $100$k transitions \\
target demonstrations             & $20$ \\
warm-up episodes (counted against the budget) & $30$ \\
batch size                        & $256$ \\
learning rate                     & $3{\times}10^{-4}$ \\
discount $\gamma$                 & $0.99$ \\
critics (clipped double-$Q$)      & $2$ \\
critic updates per environment step & $5$ \\
actor update frequency            & every $2$ critic updates \\
$n$-step return                   & $3$ \\
exploration noise std             & $0.1$ \\
demonstration fraction per batch  & $0.5$ (frozen demonstration replay) \\
target-network rate $\tau$        & $0.01$ \\
\midrule
\multicolumn{2}{l}{\emph{IBRL}} \\
\midrule
prior                             & source expert fine-tuned on $\mathcal{D}_t$, frozen \\
critic initialisation             & from scratch / from $Q_s$ (two variants) \\
prior-preference margin $m$       & $0$ / $0.01$ (two variants) \\
\midrule
\multicolumn{2}{l}{\emph{Dyna-IBRL}} \\
\midrule
prior                             & stage-2 policy head (shared with our method) \\
$Q(z,a)$ head initialisation      & regression onto $\hat r + \gamma \hat V$ \\
imagined transitions per round    & $20$k \\
imagination rollout length        & $1$ \\
synthetic buffer size             & $100$k \\
real fraction per batch           & $0.5$ \\
actor learning rate               & $1{\times}10^{-5}$ \\
actor updates per round           & $500$ \\
critic updates per round          & $5000$ \\
collection noise std              & $0.1$ \\
target-policy smoothing std / clip & $0.1$ / $0.3$ \\
\midrule
\multicolumn{2}{l}{\emph{RLPD}} \\
\midrule
prior policy                      & --- \\
offline buffer                    & $2000$ source episodes $+$ $20$ target demonstrations \\
sampling                          & symmetric (half offline, half online) \\
critic                            & layer-normalised, random subset selection \\
\midrule
\multicolumn{2}{l}{\emph{IQL} (offline warm start, then online)} \\
\midrule
offline data                      & $5000$ source episodes $+$ $20$ target demonstrations \\
expectile $\tau_{\text{exp}}$     & $0.9$ \\
offline critic steps              & $80$k \\
offline actor steps               & $20$k \\
policy extraction                 & reward-weighted regression \\
weighting temperature $\beta$, clip & $3.0$, $100$ \\
online phase                      & identical to the shared block above \\
\bottomrule
\end{tabular}
\end{table}

% ---------------------------------------------------------------------------
% PROVENANCE (read from the cfg.yaml of runs that are actually reported):
%   IBRL / RLPD / IQL  exps/simdist_threading_3arm/{ibrl,rlpd,iql}_seed1
%   Dyna-IBRL          exps/simdist_threading_3arm/dynaibrl_m001_seed1 (dyna block)
%   RLPD prior set     scripts/merge_prior_dataset.py (source rollouts + 20 demos)
%   IQL warm start     scripts/pretrain_iql_qagent.py --actor_mode rwr,
%                      then scripts/iql_ckpt_to_agent.py to remap the keys
%
% TODO(verify): RLPD's critic block. The shared table now says 2 critics for
% everything, which is what train_rl.py gives every arm -- but RLPD's published
% recipe uses a 10-critic ensemble with random subset selection, and the prose
% above claims "critic ensemble with random subset selection". Either the RLPD
% arm really does run 2 critics (in which case the prose overclaims and should
% say so, and the deviation from the published recipe must be stated) or the
% ensemble is configured somewhere I did not find. Resolve before submission --
% as written, the table and the prose contradict each other.
% TODO(verify): warm-up episodes. The RLPD cfg says num_warm_up_episode=2030,
% which is 2000 preloaded prior episodes plus the same 30 real warm-up rollouts
% the other arms get. The shared row of 30 is therefore correct, but only if the
% preloaded 2000 are genuinely not target interaction -- confirm.
% TODO(verify): Dyna-IBRL is listed under the shared 30-episode warm-up row by
% omission. It actually starts from the stage-2 prior with no separate warm-up
% phase; either give it an explicit row or exclude it from that row.
% TODO(verify): tau=0.01 in the shared block is read from the Dyna-IBRL config.
% The IQL OFFLINE phase uses 5e-3 (listed separately is fine), but check what
% train_rl.py uses for the IBRL/RLPD/IQL online phases before claiming it shared.
% TODO(check): docs/infra.md 13 warns that job_{coffee,square}_baselines.sh still
% runs the ibrl arm with warm-up 10 and WITHOUT count_warmup_transitions, unlike
% Threading. Either re-run those or state the per-task difference; as written
% this section claims one convention for all tasks.
% TODO(cite): confirm bibkeys hu2023ibrl / ball2023rlpd / kostrikov2022iql.
% ---------------------------------------------------------------------------
